Idea sztucznych sieci neuronowych narodziła się podczas prób modelowaniu pracy ludzkiego mózgu. Warren S. McCulloch oraz Walter Pitts w swojej pracy "A logical calculus of the ideas immanent in nervous activity" \cite{McCulloch1943} zaproponowali model sieci nauronowej, jako grafu skierowanego, gdzie węzły to matematyczne modele neuronu takie, że

\[N_{i}\left( t+1 \right) = H\left( \sum_{j=1}^{n}w_{ij}(t)N_{j}(t)-\Theta_{i}\left( t \right)\right)\]
gdzie:
\begin{itemize}
	\item N - stan neuronu, może przyjąć wartość 0 lub 1
	\item H - funkcja skokowa Heaviside'a
	\item $w_{ij}$ - waga połączenia między wektorem j oraz i.
	\item $\Theta$ - próg wyzwalania, neuron aktywuje się, gdy sygnał wejściowy go przekroczy.
\end{itemize}

Sieć złożona z takich neuronów jest nazywana dzisiaj perceptronem. Służy ona do zadań klasyfikacji binarnej. Potrafi określić przynależność do jednej z klas na podstawie wybranych cech. 


Jak zostało wcześniej wspomniane ten model został stworzony w określonym celu, opisaniu działania biologicznego neuronu. Można uogólnić ten model, pozwalając stworzyć bardziej uniwersalny model matematyczny.

\[ y=f\left( \sum_{i=1}^{m}w_{i}(t)x_{i}(t)+b\right) \]
gdzie:
\begin{itemize}
	\item $w_{i}$ – waga
	\item $x_{i}$ – wejście, 
	\item f – funkcja aktywacji, 
	\item y – wyjście, 
	\item m – liczba wejść,
	\item b – stała wartość (ang. bias).
\end{itemize}

W przypadku neuronu McCollocha-Pilltsa, funkcja aktywacji przyjęła postać funkcji skokowej Heaviside'a. Dzięki takiemu przedstawieniu działania neuronu, możemy zastosować dowolną funkcję aktywacji.


